% =============================================================================
% PHYSICAL AI STACK DIAGRAM (VOLUME IV)
% =============================================================================
% Renders the Volume IV margin stack. The interface mirrors Volume 1's
% \mlsysstack, Volume 2's \mlfleetstack and Volume 3's \mlagentstack: each
% argument is an intensity 0-100 for one tier, so call sites look the same
% across the series.
%
% Usage:
%   \mlphysicalstack{<governance>}{<brain>}{<nervous>}{<body>}
%   \mlphysicalstackfull{<governance>}{<brain>}{<nervous>}{<body>}
%
% Stack (top to bottom), defined by control rate rather than by dependency:
%   Governance       0.1-1 Hz     safety cases, verification, shared autonomy
%   Brain            1-50 Hz      perception, memory, intent, planning
%   ---- proposal boundary: above may only propose, below may command ----
%   Nervous System   1000 Hz      reflex enforcement, silicon placement
%   Body             continuous   matter, inertia, thermodynamics
%
% DESIGN NOTES (these are deliberate; please do not "tidy" them back)
%
% Binary, not a ramp. There are exactly two states: the active tier is solid
% accent with white type, every other tier is one flat light tint. A value
% ramp does not survive dot gain on uncoated stock, a grayscale reprint, or a
% photocopied course pack, and it also makes the column read as a gradient the
% eye tries to interpret as data rather than as a selection.
%
% Tints are restricted to 0, 22 and 90. Nothing in the 31-84 band carries a
% legible text colour: white on a 55 tint is 2.92:1 and the accent on it is
% 3.30:1. Intensities are therefore thresholded rather than interpolated; any
% value >= 50 means active.
%
% Boxed, like all three siblings, but with a real gap between tiers rather
% than a flush column. A flush column asserts that each tier rests on the one
% beneath it, which is true of Volume 1's dependency order and false here:
% these are rate bands with an authority partition, not a support hierarchy.
%
% The gap between Brain and Nervous System is deliberately several times the
% others. That is the proposal boundary, encoded by proximity rather than by a
% drawn rule, which costs no ink and needs no legend.
%
% Labels are set in real Helvetica through \usefont. A bare \sffamily is not
% usable under this font stack: with fontspec on LuaLaTeX the default encoding
% is TU, phv has no TU font, and NFSS silently substitutes a serif. The rung is
% the full 3.05cm, which fits every tier name comfortably; "Nervous System",
% the longest, measures 2.02cm in Helvetica bold at 8pt.
%
% No per-rung icons. Volume 3 dropped them and was right to: they are
% indistinguishable at 1.25in, carry nothing the label does not, and compete
% with the one mark that matters.
% =============================================================================

% --- Palette -----------------------------------------------------------------
% Solid inks rather than screened tints. A screened rule under 1pt puts about
% one screen row along its length and prints dotted or drops out entirely.
\definecolor{physctxtint}{HTML}{D3DBD7}  % inactive tier fill, solid ink
\definecolor{physctxbar}{HTML}{9AA5A0}   % rule used by the full variant
\definecolor{physctxink}{HTML}{5A6560}   % inactive tier label
\providecolor{emeraldpine}{HTML}{1A4D3E} % volume accent, if not already set

\newcommand{\physlabelfont}{\usefont{T1}{phv}{b}{n}\fontsize{8}{9}\selectfont}
\newcommand{\physratefont}{\usefont{T1}{phv}{m}{n}\fontsize{7}{8}\selectfont}

% --- One tier ----------------------------------------------------------------
% #1 y-position  #2 tier name  #3 rate  #4 intensity (0-100; >=50 is active)
\newcommand{\physrow}[4]{%
  \ifnum#4<50 \def\physbg{physctxtint}\def\phystx{physctxink}%
  \else       \def\physbg{emeraldpine}\def\phystx{white}\fi
  \fill[\physbg] (0,#1) rectangle (3.05,#1+1.02);
  \node[anchor=west,text=\phystx,font=\physlabelfont] at (0.12,#1+0.64) {#2};
  \node[anchor=west,text=\phystx,font=\physratefont]  at (0.12,#1+0.33) {#3};}

% --- Margin stack ------------------------------------------------------------
% Tier tops are 3.30 / 2.20 / 0.90 / -0.20 against a 1.02 tier height, so the
% ordinary gap is 0.08 and the Brain-to-Nervous gap is 0.28. That wider gap is
% the proposal boundary, encoded by proximity rather than by a drawn rule.
\newcommand{\mlphysicalstack}[4]{%
  \begin{tikzpicture}
    \physrow{3.30}{Governance}{0.1--1 Hz}{#1}
    \physrow{2.20}{Brain}{1--50 Hz}{#2}
    \physrow{0.90}{Nervous System}{1000 Hz}{#3}
    \physrow{-0.20}{Body}{continuous}{#4}
  \end{tikzpicture}%
}

% --- Full-size stack for the About page --------------------------------------
% Carries the subtitles and names the proposal boundary once, in the one place
% the reader is being introduced to the architecture rather than located in it.
\newcommand{\physrowfull}[5]{%
  \ifnum#5<50
    \fill[physctxbar] (0.22,#1+0.10) rectangle (0.34,#1+1.30);
    \def\pfink{physctxink}%
  \else
    \fill[emeraldpine] (0,#1) rectangle (0.46,#1+1.40);
    \def\pfink{emeraldpine}%
  \fi
  \node[anchor=west,text=\pfink,
        font=\usefont{T1}{phv}{b}{n}\fontsize{10.5}{12}\selectfont]
        at (0.76,#1+1.02) {#2};
  \node[anchor=west,text=\pfink,
        font=\usefont{T1}{phv}{m}{n}\fontsize{8.5}{10}\selectfont]
        at (0.76,#1+0.70) {#3};
  \node[anchor=west,text=\pfink,
        font=\usefont{T1}{phv}{m}{n}\fontsize{8}{9.5}\selectfont]
        at (0.76,#1+0.38) {#4};}

\newcommand{\mlphysicalstackfull}[4]{%
  \begin{tikzpicture}
    \physrowfull{4.80}{Governance}{0.1--1 Hz}%
      {safety cases, verification, shared autonomy}{#1}
    \physrowfull{3.20}{Brain}{1--50 Hz}%
      {perception, memory, intent, planning}{#2}
    \draw[physctxbar,line width=0.9pt,dash pattern=on 2.2pt off 1.8pt]
          (0,3.02) -- (8.6,3.02);
    \node[anchor=west,text=physctxink,
          font=\usefont{T1}{phv}{m}{n}\fontsize{7.5}{9}\selectfont]
          at (0,2.76) {above this line a tier may only propose; below it a tier may command};
    \physrowfull{1.20}{Nervous System}{1000 Hz}%
      {reflex enforcement, silicon placement}{#3}
    \physrowfull{-0.40}{Body}{continuous}%
      {matter, inertia, thermodynamics}{#4}
  \end{tikzpicture}%
}
